% !TeX root = ../papers/nc-p3.tex
\section{可逆な交差ブロックを持つ二つの疎長方形の全補完}
\label{sec:parallel-invertible-rectangles}
\begin{theorem}\label{thm:parallel-invertible-rectangles}
平行二重色の表で76--85または112--147と同じ薄い部分の台を持つ窓を取る。型をそれぞれ76型、112型と呼ぶ。
厚い部分には追加の非零条件を課さない。各型で、位置別対角変更を除いた二切断階数の解は
一つの有理的な主成分と、$t=1,-1$ 上の二つの成分で尽くされる。その係数・交わりは以下の式で与えられる。
主成分の幾何学的一般点は標準Koszul左右AS正則、反対階数16、最小周期4で、
無印の像の次元はそれぞれ3,4である。$t=\pm1$ の二つの特別成分では、後述の例外集合を除く全点が標準正則である。
\end{theorem}

\begin{proof}[薄い部分の規格化]
前節の16台のうち88--90を除いた全てがこの二形のどちらかに含まれる。
薄い色2,3の四長方形で非零成分の二次小行列式を解く。指数行列のSmith非零対角は全て1で、
位置別倍率を除くと一つの可逆な変数 $t$ が残る。検算で用いる順序の規格化は次の通りである。
薄い第一切断の二つの行列を順に、列ベクトルと行ベクトルの外積で書けば
\[
\begin{array}{c|cc}
 & M_2 & M_3\\\hline
76 &(t,0,1,1)^t(1,1,1,0)&(1,1,1,0)^t(1,0,1,1)\\
112 &(t,1,1,1)^t(1,1,1,0)&(t,1,t,0)^t(1,1,1/t,1/t).
\end{array}
\]
薄い第二切断が残る厚いブロックの右上 $U$、左下 $V$ を決定する。
各二行行列を四成分の行順のリストで書き、切断・色を $(0,0),(0,1),(1,0),(1,1)$ の順とすると
\[
\begin{array}{c|l|l}
 &U&V\\\hline
76 &(t,0,0,1);(1,t,1,0);(1,0,0,1);(0,1,1,1)
   &(1,t,1,0);(1,0,0,t);(t,1,1,0);(1,0,0,1)\\
112&(t,0,1,1);(t,t,1,0);(1,0,t,1);(0,1,1,t)
   &(1,1,1,0);(1,1,0,1);(t,1/t,1,0);(1/t,0,1,1).
\end{array}
\]
いずれも可逆であり、各厚い行列が階数2となる条件は左上 $H$ が可逆で右下が $V H^{-1}U$ となることである。
\end{proof}

\begin{proof}[行列式比による全消去]
左上の四行列は順に
\[
 (x_0,x_1,x_6,x_7),\ (x_2,x_3,x_4,x_5),\
 (x_0,x_4,x_3,x_7),\ (x_5,x_1,x_6,x_2)
\]
となる。それらを $H_0,\ldots,H_3$ とし、$\det H_0/\det H_i$ を $i=1,2,3$ の順に
$a,b,c$ と置く。第一切断の右下を $V H^{-1}U$ で定め、第二切断が同じ語に要求する右下と比較する。
共通の $\det H_0$ を払うと、$V\operatorname{adj}(H_i)U$ に順に $1,a,b,c$ を掛けた成分の比較となり、$x_0,\ldots,x_7$ に関する八本の斉次一次式が得られる。
係数行列の行列式は $(at^2-bc)^4$（112型ではさらに $t^{-4}$ 倍）である。
従って $a=bc/t^2$。$A=x_2,B=x_4,C=x_5,D=x_6$ と置く。
両型とも $x_1=bB$ で、残る三つは
\[
\begin{array}{c|ll}
 &76 &112 \\\hline
t x_0 & cA+(t^2-c)B & cA+t(t^2-c)B\\
t x_7 & cC & (t^2-c)B+cC\\
b x_3 &t^2D+(b-1)cC &
 t^2D+(b-1)cA+t(b-1)cC+t(t^2(b-2)-2bc+b+2c)B .
\end{array}
\]
ここまでに割るのはトーラス単元だけである。$\delta=t^2+bc-b-c$ と書く。
残る一次式は $\delta B=0$ と、76型なら $\delta C=0$、112型なら $\delta(A-2tB+tC)=0$ となる。
$\delta\ne0$ のとき前者は $\det H_0=0$ を与える。後者では $B=0,A=-tC$ と代入すると
三つの行列式比の条件は各々単元倍の $C^2(b-c),C^2(b-1),C^2(t^2-c)$ になる。
$C=0$ なら再び行列式が零、$C\ne0$ なら $\delta=0$ に反する。よって必ず
\[
 t^2=b+c-bc .
\]
この曲面上、三つの行列式比の条件は、共通単元と符号を除いて
$(b-c)F,c(b-1)F,b(c-1)F$ である。ただし
\[
 F_{76}=t^2BD+cC(A-B),\qquad
 F_{112}=-t^2BD+cA(B-C)+tb(1-c)B^2+tcBC .
\]
これで $tbc\det H_0\ne0$ の下、全ての係数条件は
$F=0$ または $b=c=1,\ t=\pm1$ となる。逆に三つの行列式比が正しければ他の三行列式も非零で、
右下の二表示が一致し全ての階数を得る。

薄い部分の対角安定群の作用は112型では $A,B,C,D$ の共通倍率であり、
76型では $(A,B)$ と $(C,D)$ の独立倍率である。これは薄い語の位置別指数行列の整数核で直接分かる。
112型のSmith因子2による残余符号も全厚い座標に自明に作用する。
$F=0$ 上で $B=0$ なら $AC=0$ で $\det H_0=0$ となる。従って $B=1$ にでき、$D$ は一次式から一意である。
76型では $C\ne0$ も強制されるので $C=1$ とできる。基底曲面は $b\ne1$ で $c=(t^2-b)/(1-b)$、
$b=1$ なら $t=\pm1$ で、曲面全体は既約である。
従って主成分の次元は3,4で、各特別成分は $AC-BD\ne0$ の開部分を持つ
$\mathbb P^1\times\mathbb P^1$, $\mathbb P^3$ である。それぞれの主成分との交わりは同じ $F=0$ である。
\path{computations/verify_parallel_invertible_rectangles.py} がSmith変換、八行列式、残る一次式と三つの二次式を検算する。
\end{proof}

\begin{proof}[四つの特別成分の全点]
$b=c=1,t=\pm1$ とする。$t=1$ では第四位置の文字に1を加えると偶二位置の交換と奇二位置の交換を個別に保つ。
$t=-1$ でも先に四元群捻りを戻せば同じである。実際、さらに第三位置に符号 $(1,1,-1,1)$、
第四位置の平行移動前に76型で $(1,-1,1,1)$、112型で $(-1,1,-1,1)$ を掛ければよい。
この両対称形の隣接六平面には純テンソルがない。混合色の二つの関係を係数 $u,v$ で取ると、
二つの混合二行ブロックの行列式は $u^2,v^2$ の独立な一次式であり、純関係なら $u=v=0$ となる。
残る対角二ブロックはどちらへの射影も同型であるため、片方が零なら他方も零である。

この両対称表示で、全ての非零片側評価の階数が3以上であることが標準正則性の必要十分条件である。
実際、総色1のため各側の平方四成分には少なくとも二次元の対角反対核がある。
次元3以上なら二つの対角成分を零にでき、可約な二次式を持つ。対角鉛筆が二次元なら、
四列のどれかが零か二列が比例すれば同じく可約な元を持ち、そうでなければ滑らかな四次完全交叉になる。
反対切断の混合色2,3は薄い部分だけで決まり常に階数2である。従って両側で対角核が二次元なら全反対階数は8である。
可約鉛筆では補題\ref{lem:bisymmetric-reducible-quick-test}、両側滑らかなら
定理\ref{thm:smooth-apolar-complete}を適用する。必要性は両対称窓の片側評価障害による。

実際に階数2以下の評価があるのは、$AC-BD\ne0$ の中で次に挙げる各行の条件の和集合に限り、各条件で必ず存在する。
\[
\begin{array}{c|l}
76 &(D=0,\ A=B),\quad(C=D,\ A=0)\\
112 &(B=tD,\ A=D),\quad(C=tD,\ B=tD),\quad(C=tD,\ A=0),\quad(C=0,\ A=D).
\end{array}
\]
これを確かめるには、始点0,1において厚い関係に $\binom{H}{V}$ の二列、
薄い関係にその非零列を一つ取る。$t=-1$ では捻りを戻す行符号を加える。
文字対の一方だけの評価では厚い部分が相手の一対を張り、薄い関係の少なくとも一つが他対に非零となる。
両対で非零なら、階数2以下の像は両対に同型に射影するグラフである。色0関係に限ると両対の係数射影が可逆なので、
評価座標の非零個数が各対の間で一致する。従って調べる台は各対一文字ずつの四通りと全四文字だけである。
各台で最初の座標を1にし、残る非零座標と $AC-BD$ を可逆化する。六行四列の評価行列の三次小行列式の消去根基を求めると、
表示の各一次条件対または単位イデアルを得る。評価点全体は射影的で、階数の下がる射影点の存在は基底上の閉条件である。
したがって各台の像の閉包の和集合を取れば、特殊値も含めて正確な不正則集合になる。
\path{computations/verify_parallel_vertical_regularity.py} は40個の消去根基で両方向のイデアル包含を確認する。
\end{proof}

\begin{proof}[主成分の一般点]
各型の $t=b=c=1$, $(A,B,C)=(2,1,3)$, $F=0$ の点は上の正則判定を通る
（$D$ は順に $-3,-1$）。主成分はこの点を含み、二切断・一文字切断の階数が保たれる。
定理\ref{thm:klein-elliptic-completion}で用いた次数別の開条件を適用する。
すなわち正則点の近くで縮約による三次・四次交叉の下界と上半連続性により有限Frobenius双対が延び、
各内部次数のbar複体の非対角ホモロジーの消滅は一般点へ移る。よって幾何学的一般点が全次数で標準ASとなる。

$t=2,b=3,c=-1/2,(A,B,C)=(7,1,5)$ とし $F=0$ で $D$ を選ぶ。両型で反対四ブロックが可逆、
四位置の無限小線型作用の行列が階数61である。後者の核は自明な倍率だけなので一般の射影安定群は有限である。
有限群への四元群の埋込みは有限個で、各作用で文字が異なるため、印を忘れても像の次元は減らない。
さらに四次交代収縮の勾配 $G$ と値 $D_4$ を使い、各始点 $j$ の隣接切断で $L_j=W_jG_j^t/D_4$ とする。
第一文字を残す部分traceで $\operatorname{tr}((\operatorname{tr}_{j+1} L_j^2)^2)$ はこの有理点の四始点で相異なる。
よって一般点に始点を動かす周期1,2はなく、最小周期は4である。これらの行列計算も前記の補完検算に含まれる。
\end{proof}
